% GENERATED FILE. Edit canonical lessons, then run npm run generate:books.
% canonical-source-hash: fnv1a64:000d4b540fa7f49a
% canonical-lessons: ES-C361-carne, ES-C361-pescado, ES-C361-leche, ES-C361-sal, ES-C361-aceite

\chapter{What Is On The Table}
\label{ch:la-carne-y-la-leche}
\hlchaptermodality{361}

\begin{chapteropening}
\textbf{By the end of this chapter:} \emph{I can say la carne, el pescado, la leche, la sal and el aceite, and ask for the things a plain meal is made of.}

Complete ES-C361-aceite: name two things you cook and three things you cook them with.
\end{chapteropening}

\section[la carne]{la carne — the flesh that named a party before a fast}

\label{lesson:ES-C361-carne}



\begin{quote}
Say \emph{the trousers}, then \emph{the button}. (\emph{El pantalón}; \emph{el botón}.)
\end{quote}

\subsection*{What to know first}
\begin{itemize}
\raggedright
  \item la comida — the meal this word ends up inside.
  \item el plato — what the meat arrives on.
\end{itemize}

\begin{sounds}
\begin{itemize}
\raggedright
  \item \emph{KAR-ne}, two beats with the weight on the first. The \emph{r} is a single tap of the tongue, not a roll. $\to$ pronunciation reference
\end{itemize}
\end{sounds}

\begin{cousinweb}
\textbf{la carne} — \textquotedblleft{}meat.\textquotedblright{}

Latin \textbf{carō} had the stem \textbf{carn-}, and it meant flesh: the stuff of a body, on the animal or on the person. Spanish took it plain, and \emph{la carne} is what a butcher sells and what arrives for dinner. No shudder attached.

English inherited nothing here and borrowed everything, which is why the same root arrives in English wearing solemn clothes:

\begin{itemize}
\raggedright
  \item \textbf{carnal}, of the flesh rather than the spirit;
  \item \textbf{incarnate}, made flesh, out of the Latin of the Church;
  \item \textbf{carnage}, through French and Italian — flesh in quantity.
\end{itemize}

The generous one is \textbf{carnival}. Italian \emph{carnevale} goes back to an older \emph{carnelevare}, \textquotedblleft{}the removing of meat\textquotedblright{} — the name of the hinge in the year where meat stops and Lent begins. The party is named after the fast that follows it.

You will also hear a tidier account: \emph{carne, vale!}, \textquotedblleft{}meat, farewell!\textquotedblright{}, as though the word were a parting wave at a plate. It scans beautifully and it is not where the word came from. It is a reworking by speakers who heard \emph{carnevale}, recognised \emph{carne} inside it, and supplied a story for the rest.

Worth hearing too: the register gap. English \emph{carnal} is a heavy word, kept for sermons and courtrooms. Spanish \emph{carne} is the lightest one there is. It goes on the shopping list.
\end{cousinweb}

\begin{grammarlens}[title={What you've built}]
You can name what is on a plate with the same root that named a street party in February: \emph{la carne}.
\end{grammarlens}

\subsection*{Your turn}
Say these aloud:

\begin{itemize}
\raggedright
  \item \emph{la carne}
  \item \emph{la carne y el pan}
  \item \emph{un plato de carne}
\end{itemize}

\begin{tcolorbox}[breakable,colback=teal!4,colframe=teal!35!black,title={Before you move on}]
What did \emph{carnelevare} mean? (The removing of meat, before Lent.) And is \emph{carne, vale} the real story? (No — that is a later guess, made by ear.)
\end{tcolorbox}
\section[el pescado]{el pescado — the fish that has already been caught}

\label{lesson:ES-C361-pescado}



\begin{quote}
Say \emph{the button}, then \emph{the meat}. (\emph{El botón}; \emph{la carne}.)
\end{quote}

\subsection*{What to know first}
\begin{itemize}
\raggedright
  \item el agua — where the fish was before it was caught.
  \item la olla — where it goes afterwards.
\end{itemize}

\begin{sounds}
\begin{itemize}
\raggedright
  \item \emph{pes-KA-do}, three beats with the weight on the middle one. The final \emph{d} is soft, close to the \emph{th} of \emph{the}. $\to$ pronunciation reference
\end{itemize}
\end{sounds}

\begin{cousinweb}
\textbf{el pescado} — \textquotedblleft{}fish,\textquotedblright{} the kind you eat.

Latin \textbf{piscis} was a fish. From it came the verb \emph{piscārī}, to fish, and from that verb its past participle \textbf{piscātus}, \textquotedblleft{}having been fished\textquotedblright{} — the catch. Spanish turned that participle into a noun. So \emph{el pescado} is not a fish. It is \textbf{the caught one}.

That is why Spanish is holding two words. A fish moving through water is \emph{un pez}. The same creature, netted and in a pan, is \emph{el pescado}. The grammar of a finished action is doing the work, and you can hear the \emph{-ado} on the end of it, the same ending that sits on the end of \emph{hablado} and \emph{comido}.

English took the learned half straight from Latin: \textbf{Pisces} in the sky, \textbf{piscine} for anything fishy, \textbf{piscatorial} for the business of fishing. But English also owns the word by inheritance. Latin \emph{pater} is English \emph{father}; Latin \emph{pēs} is \emph{foot}; Latin \emph{plēnus} is \emph{full}. \emph{Piscis} and \emph{fish} are that same pairing, one more time.

A caution on how far back to push it. Call it the p and f correspondence and stop there — the fish-word turns up only in the western languages, and one proposal is that they all took it from some older neighbour rather than inheriting it.

One letter of danger: \emph{pescado} is fish, \emph{pecado} is sin. \emph{Pescador} fishes, \emph{pecador} sins.
\end{cousinweb}

\begin{grammarlens}[title={What you've built}]
You can order fish with a word that means the catch, and hear a finished action inside it: \emph{el pescado}.
\end{grammarlens}

\subsection*{Your turn}
Say these aloud:

\begin{itemize}
\raggedright
  \item \emph{el pescado}
  \item \emph{el pescado y el arroz}
  \item \emph{un plato de pescado}
\end{itemize}

\begin{tcolorbox}[breakable,colback=teal!4,colframe=teal!35!black,title={Before you move on}]
What does the \emph{-ado} of \emph{pescado} tell you? (That it has been caught — it is a past participle.) And what does \emph{pecado} mean? (Sin — an unrelated word, one letter away.)
\end{tcolorbox}
\section[la leche]{la leche — the milk hiding inside a lettuce}

\label{lesson:ES-C361-leche}



\begin{quote}
Say \emph{the meat}, then \emph{the fish}. (\emph{La carne}; \emph{el pescado}.)
\end{quote}

\subsection*{What to know first}
\begin{itemize}
\raggedright
  \item la vaca — where the milk starts.
  \item el queso — what it turns into.
\end{itemize}

\begin{sounds}
\begin{itemize}
\raggedright
  \item \emph{LE-che}, two beats with the weight on the first. The \emph{ch} is one sound, the \emph{ch} of \emph{church}. $\to$ pronunciation reference
\end{itemize}
\end{sounds}

\begin{cousinweb}
\textbf{la leche} — \textquotedblleft{}milk.\textquotedblright{}

Latin \textbf{lac} had the stem \textbf{lact-}, and the spoken language flattened it to \emph{lacte}, which wore down into \emph{leche}. Chemists later went back to the Latin stem and cut it into \textbf{lactose}, \textbf{lactic} and \textbf{lactate} — modern coinages from an old root. Italian \textbf{latte} needs no explanation at all: it is \emph{leche}, in another accent.

Now the good one. A \textbf{lettuce} is a \emph{lactūca}, \textquotedblleft{}the milky one,\textquotedblright{} because a cut lettuce stem bleeds a pale sap. English got the word through Old French \emph{laitues}. Spanish never had to borrow it — \emph{lechuga} came down the ordinary way, so Spanish keeps the milk and the salad plant as an audible pair, no translation required.

\textbf{Milky Way} is a translation of Latin \emph{via lactea}, which was itself a translation of a Greek phrase — the same picture handed along in three languages.

And a hedge on the Greek. \textbf{Galaxy} comes from \emph{galaxías}, \textquotedblleft{}milky\textquotedblright{}, from Greek \emph{gála}, stem \emph{galakt-}. That much is solid. What is not solid is treating \emph{gála} and \emph{lac} as a tidy pair of cousins: the sounds do not line up the way inherited words normally do, and the usual account now is that Greek and Latin each picked their milk-word up from somewhere else that has been lost. Two words that look like damaged copies of each other, with no surviving original.
\end{cousinweb}

\begin{grammarlens}[title={What you've built}]
You can name milk, and hear the same word inside a salad leaf: \emph{la leche}.
\end{grammarlens}

\subsection*{Your turn}
Say these aloud:

\begin{itemize}
\raggedright
  \item \emph{la leche}
  \item \emph{la leche y el queso}
  \item \emph{un vaso de leche}
\end{itemize}

\begin{tcolorbox}[breakable,colback=teal!4,colframe=teal!35!black,title={Before you move on}]
Why is a lettuce named for milk? (A cut stem bleeds a milky sap.) And are Greek \emph{gála} and Latin \emph{lac} a clean pair of cousins? (No — they look like damaged copies of each other, and neither seems inherited.)
\end{tcolorbox}
\section[la sal]{la sal — one short word, and three roots that only sound alike}

\label{lesson:ES-C361-sal}



\begin{quote}
Say \emph{the fish}, then \emph{the milk}. (\emph{El pescado}; \emph{la leche}.)
\end{quote}

\subsection*{What to know first}
\begin{itemize}
\raggedright
  \item la salud — the sal- word that is not salt.
  \item la sopa — the first thing you reach for salt to fix.
\end{itemize}

\begin{sounds}
\begin{itemize}
\raggedright
  \item \emph{sal}, one beat. The \emph{s} stays clear and unbuzzed, and the \emph{l} is light, made at the tip of the tongue. $\to$ pronunciation reference
\end{itemize}
\end{sounds}

\begin{cousinweb}
\textbf{la sal} — \textquotedblleft{}salt.\textquotedblright{}

Latin \textbf{sal, salis}. Its family is a kitchen: \textbf{salad} is \emph{herba salāta}, salted greens; \textbf{salsa} and \textbf{sauce} are one Latin word, \emph{salsa}, \textquotedblleft{}salted\textquotedblright{}, arriving in English by two different doors; \textbf{sausage} is \emph{salsīcia}; and \textbf{saline} and \textbf{salami} are the same root again.

\textbf{Salary} belongs here too — \emph{salārium} is built on \emph{sal}, and that much is real. What is not real is the story stapled to it. No Roman source says soldiers were paid in salt. Pliny remarks only that salt was valuable enough for the word to be named after it; the salt-money version traces to an eighteenth-century dictionary. It is the same species of invention as \emph{carne, vale!}.

Now the part that matters for the ear. Spanish has other words that open with \emph{sal-}, and they are not a crowd of separate accidents. There are three roots in total, counting salt.

\begin{itemize}
\raggedright
  \item \textbf{Whole, safe.} \emph{Salūs} gives \emph{salud}; \emph{salūtāre}, built directly on \emph{salūs}, gives \emph{saludo}; \emph{salvus} gives \emph{salvo}. Those three Spanish words are \textbf{one root}, not three. English drew \emph{salute}, \emph{salvage}, \emph{salvation} and \emph{save} from it, and by way of Greek \emph{hólos} the \emph{holo-} of \emph{holograph}.
  \item \textbf{Leap.} \emph{Salīre} meant to jump. English took \emph{salient}, \emph{sally}, \emph{assail}, \emph{assault}, \emph{result}, \emph{insult} and \emph{somersault} from it. Spanish \emph{salir}, to go out, is a leap that calmed down into a doorway.
\end{itemize}

So: one root seasons the food, one wishes you well, one jumps.
\end{cousinweb}

\begin{grammarlens}[title={What you've built}]
You can salt a soup, and tell three \emph{sal-} families apart by ear: \emph{la sal}.
\end{grammarlens}

\subsection*{Your turn}
Say these aloud:

\begin{itemize}
\raggedright
  \item \emph{la sal}
  \item \emph{la sal y el pan}
  \item \emph{un poco de sal}
\end{itemize}

\begin{tcolorbox}[breakable,colback=teal!4,colframe=teal!35!black,title={Before you move on}]
How many roots hide behind the \emph{sal-} words? (Three — salt, whole-or-safe, and leap.) And which Spanish words share the whole-or-safe one? (\emph{Salud}, \emph{saludo} and \emph{salvo}, all from the same root.)
\end{tcolorbox}
\section[el aceite]{el aceite — an Arabic oil and a Latin olive in one phrase}

\label{lesson:ES-C361-aceite}



\begin{quote}
Say \emph{the milk}, then \emph{the salt}. (\emph{La leche}; \emph{la sal}.)
\end{quote}

\subsection*{What to know first}
\begin{itemize}
\raggedright
  \item la sartén — what you pour it into.
  \item la cebolla — what goes in after it.
\end{itemize}

\begin{sounds}
\begin{itemize}
\raggedright
  \item \emph{a-SEI-te}, three beats with the weight on the middle one. In much of Spain that \emph{c} is the soft \emph{th} of \emph{think}; across the Americas it is a clear \emph{s}. $\to$ pronunciation reference
\end{itemize}
\end{sounds}

\begin{cousinweb}
\textbf{el aceite} — \textquotedblleft{}oil.\textquotedblright{}

Arabic \textbf{az-zayt} means \textquotedblleft{}the oil.\textquotedblright{} Spanish borrowed the whole phrase, article and all, and never took the article back off. The Arabic \emph{al-} had already softened to \emph{az-} alongside that \emph{z}, so what crossed over was \emph{azzayt}, which settled into \emph{aceite}. The olive, \emph{az-zaytūna}, came the same way and became \emph{aceituna}. English did this too, with \textbf{alcohol}, \textbf{algebra} and \textbf{alchemy} — each of them carrying an Arabic \emph{al-} that means nothing at all in English now.

Latin had a perfectly good word of its own: \textbf{oleum}, oil, and \textbf{olīva}, the tree and its fruit. That is the source of English \textbf{oil} (through Old French), of \textbf{olive}, of \textbf{petroleum}, \textquotedblleft{}rock oil\textquotedblright{}, and of \textbf{linoleum}, a nineteenth- century coinage from \emph{linum}, flax, plus \emph{oleum}. Spanish kept \emph{óleo} for the painter's oil and \emph{el olivo} for the tree.

Here is the part to keep. Neither word is home-grown. Arabic \emph{zayt} is Semitic — Hebrew \emph{zayit} is the same word — and it has no connection to \emph{oleum}. And \emph{oleum} and \emph{olīva} were themselves lifted from Greek \emph{élaion} and \emph{elaíā}, which are generally taken to come from some older Mediterranean language that left nothing else behind. So the phrase \emph{aceite de oliva} stands two borrowings next to each other, a thousand years apart, naming one plant.

One warning for the kitchen and the garage: \emph{aceite} is oil of every kind, cooking or engine. The stuff from the ground is \emph{petróleo}.
\end{cousinweb}

\begin{grammarlens}[title={What you've built}]
You can name the oil in a pan with a word Spain took from Arabic, and the olive in it with one Rome took from Greek: \emph{el aceite}.
\end{grammarlens}

\subsection*{Your turn}
Say these aloud:

\begin{itemize}
\raggedright
  \item \emph{el aceite}
  \item \emph{el aceite y la sal}
  \item \emph{un poco de aceite}
\end{itemize}

\begin{tcolorbox}[breakable,colback=teal!4,colframe=teal!35!black,title={Before you move on}]
What is the \emph{a-} at the front of \emph{aceite}? (The Arabic article \emph{al-}, already softened to \emph{az-} before the \emph{z}.) And where did Latin get \emph{oleum}? (Borrowed from Greek — so both words for the plant are loans.)
\end{tcolorbox}
